\section{De las semillas APP al registro de mezcla y masa}
\label{sec:semillas-registro-mezcla-autonoma}
\index{semilla!factorización suma--producto}
\index{mezcla!procedencia en las semillas}

La ley física recibe un dominio que ya ha conservado la incompatibilidad de
las dos evaluaciones APP. Si
\(\mathcal S_+\) y \(\mathcal S_\times\) denotan las semillas orientadas de
las hojas aditiva y multiplicativa, entonces
\begin{equation}
 \Omega_{\rm seed}=\mathcal S_+\times\mathcal S_\times,
 \qquad
 |\Omega_{\rm seed}|=324^2=104,976.
 \label{eq:dominio-semillas-md-autonoma}
\end{equation}
La factorización exacta del emisor produce
\begin{equation}
 104,976
 \longrightarrow18\times26=468\ U_6
 \longrightarrow243\ w_6\subset\mathbb F_3^6,
 \qquad |\mathbb F_3^6|=729.
 \label{eq:embudo-semillas-md-autonoma}
\end{equation}
Los cuatro cardinales cuentan objetos distintos: condiciones iniciales de
dos hojas, emisiones decimales, palabras ternarias publicadas y ambiente
algebraico. Cada semilla conserva, además de su emisión, multiplicidad,
diagonal, hoja, órbita, firma de emisión, residuos módulo tres y módulo
nueve, giros, cambios de hoja, extremo, desplazamientos y conteos
cardinales. Esos datos constituyen el vocabulario del registro posterior.

La factorización organiza seis microfibras de cardinal \(1008\), seis planos
de Fano de fase, treinta y seis fibras de cardinal \(28\) y la banda
\(1944=27\cdot72\), diagonalizada por \((\mathbb Z/3)^3\). Al promediar la
palabra nonádica del TPK, los proyectores condicionales de las dos hojas
definen el núcleo canónico
\begin{equation}
 K_{\rm cat}=\frac13(I+E_++E_\times),
 \qquad
 \operatorname{Spec}(K_{\rm cat})
 =\left\{1^{[1]},
 \left(\frac23\right)^{[42]},
 \left(\frac13\right)^{[425]}\right\}.
 \label{eq:kernel-canonico-semillas-md-autonoma}
\end{equation}
La fase TPK selecciona la renovación activa y conserva las nueve bandas
espaciales; el promedio~\eqref{eq:kernel-canonico-semillas-md-autonoma}
publica la sombra estocástica del calendario completo.

\subsection{Bidegraduación de las semillas y \(6\) planos de incidencia}

En cada una de las seis microfibras de cardinalidad \(1008\), el grafo
bipartito conserva \(54\) vértices aditivos y \(56\) multiplicativos. Los
primeros se descomponen en \(36\) vértices de grado \(24\) y \(18\) de
grado \(8\); los segundos tienen grado \(18\). El recuento de incidencias
coincide por ambas hojas:
\begin{equation}
 1008=36\cdot24+18\cdot8=56\cdot18.
 \label{eq:bidegraduacion-microfibra-masas-autonoma}
\end{equation}
Esta igualdad fija la bidegraduación del emisor antes de proyectarlo sobre
las familias de ruta.

Las seis palabras
\begin{equation}
 001112,\quad010211,\quad100121,\quad112001,\quad121100,\quad211010
 \label{eq:seis-palabras-microfibra-masas-autonoma}
\end{equation}
forman una órbita de \(D_3\simeq S_3\) y poseen el perfil
\(432+4\cdot144=1008\). El cociente \(1008=36\cdot28\) organiza siete
bloques
\[
 Q=\{B_3,B_6,B_9,P_1,P_2,P_3,f\},
\]
donde \((B_3,B_6,B_9)\) son las tétradas transversales y
\((P_1,P_2,P_3,f)\), las cuatro tétradas laterales. Para una de las seis
palabras \(a\), sea
\[
 \mathcal R_a
 :=\{(p,t)\in\mathcal S_+\times\mathcal S_\times:
       U_6(p,t)\text{ publica la palabra }a\}
\]
el conjunto de sus mil ocho rutas, y pongamos
\(R_9:=\mathbb Z/9\mathbb Z\). Fijada la carta orientada, la proyección
de fase tiene tipo
\begin{equation}
 q_\beta:\mathcal R_a\longrightarrow\binom{R_9}{2},
 \qquad
 q_\beta(p,t)
 =\{x-5\beta(d_t),x+5\beta(d_t)\}\subset\mathbb Z/9\mathbb Z,
 \qquad x=5(i_p-j_p),
 \label{eq:proyeccion-fano-semillas-masas-autonoma}
\end{equation}
con \((\beta(N),\beta(E),\beta(S),\beta(O))=(1,2,3,4)\). Su codominio es
el conjunto de los treinta y seis pares no ordenados de \(R_9\). Cada
\(X\in Q\) es una tétrada de órbitas;
su soporte de rutas y su corte sobre una fibra se definen por
\[
 \operatorname{Supp}_{\rm rutas}(X)
 :=\bigcup_{\mathcal O\in X}\mathcal O,
 \qquad
 X_b:=\operatorname{Supp}_{\rm rutas}(X)\cap q_\beta^{-1}(b),
 \qquad b\in\binom{R_9}{2}.
\]
En el canal multiplicativo, el observable fuente es
\begin{equation}
 \rho:\mathcal S_\times\longrightarrow R_9,
 \qquad
 \rho(i,j,d)=(i+1)(j+1)\pmod9.
 \label{eq:residuo-producto-fano-semillas-rev3}
\end{equation}
Si \(\pi_\times:\mathcal R_a\to\mathcal S_\times\),
\(\pi_\times(p,t)=t\), es la segunda proyección, el observable que actúa
sobre rutas es
\begin{equation}
 \rho_\times:=\rho\circ\pi_\times:\mathcal R_a\longrightarrow R_9,
 \qquad
 \rho_\times(p,t)=(i_t+1)(j_t+1)\pmod9.
 \label{eq:residuo-producto-rutas-fano-semillas-rev3}
\end{equation}
En lo que sigue, \(\rho(X)\) abrevia el multiconjunto de valores de
\(\rho_\times\) sobre \(X_b\), que es independiente de la fibra \(b\). Sus
perfiles sobre las tétradas son
\[
 \rho(B_3)=3^4,
 \qquad
 \rho(B_6)=6^4,
 \qquad
 \rho(B_9)=0^4,
\]
y, para las tres tétradas laterales,
\begin{equation}
 \begin{array}{c|c|c}
 r&\text{firma activa de }P_r&R(P_r)\\
 \hline
 1&933&\{5,7\}\\
 2&393&\{4,8\}\\
 3&339&\{1,2\}.
 \end{array}
 \label{eq:perfiles-laterales-fano-semillas-rev3}
\end{equation}
Cada residuo de \(R(P_r)\) aparece dos veces entre las cuatro rutas de
\((P_r)_b\).

\begin{theorem}[Morfismo de incidencia de las semillas y seis planos de Fano]
\label{thm:semillas-seis-fano-fase-rev3}
Para \(\delta\in R_9\), \(h\in\{3,6,0\}\) ---correspondiente,
respectivamente, a \(B_3,B_6,B_9\)--- y \(r\in\{1,2,3\}\), sea
\begin{equation}
 N_\delta(B_h,P_r)
 :=\sum_{b\in\binom{R_9}{2}}
 \#\left\{(x,y)\in(B_h)_b\times(P_r)_b:
 \rho_\times(y)-\rho_\times(x)=\delta\right\}.
 \label{eq:conteo-incidencia-fano-semillas-rev3}
\end{equation}
Entonces
\begin{equation}
 N_\delta(B_h,P_r)=
 \begin{cases}
 288,&h+\delta\in R(P_r),\\
 0,&h+\delta\notin R(P_r).
 \end{cases}
 \label{eq:conteo-288-cero-fano-semillas-rev3}
\end{equation}
Si \(\delta\in R_9^\times\), para cada \(B_h\) existe un único \(P_r\)
con conteo \(288\); por tanto se obtiene una biyección
\[
 \theta_\delta:\{B_3,B_6,B_9\}
 \xrightarrow{\sim}\{P_1,P_2,P_3\}.
\]
Escribiendo la imagen de \((B_3,B_6,B_9)\) mediante sus subíndices, las seis
biyecciones son
\begin{equation}
 \begin{array}{c|c@{\qquad}c|c}
 \delta&\theta_\delta&\delta&\theta_\delta\\
 \hline
 1&(2,1,3)&5&(2,3,1)\\
 2&(1,2,3)&7&(3,2,1)\\
 4&(1,3,2)&8&(3,1,2).
 \end{array}
 \label{eq:tabla-theta-delta-fano-semillas-rev3}
\end{equation}
Sea \(P=\{P_1,P_2,P_3\}\) y, para \(p\in P\), sea
\[
 \mu(p):=\bigl\{\{f,p\},P\setminus\{p\}\bigr\}
\]
el emparejamiento perfecto asociado sobre las cuatro tétradas laterales.
Entonces
\begin{equation}
 \mathcal D_\delta
 :=\bigl\{\{B_3,B_6,B_9\}\bigr\}
 \cup
 \bigl\{\{B,q,q'\}:B\in\{B_3,B_6,B_9\},
       \{q,q'\}\in\mu(\theta_\delta(B))\bigr\}
 \label{eq:sistema-fano-semillas-rev3}
\end{equation}
es un sistema de Steiner \(\operatorname{S}(2,3,7)\). Las seis unidades
producen los seis planos de Fano sobre estos siete bloques que contienen la
línea transversal \(\{B_3,B_6,B_9\}\).
\end{theorem}

\begin{proof}
Fijemos una fibra \(b\). En \((B_h)_b\) existen cuatro elecciones con
residuo \(h\). Si \(h+\delta\in R(P_r)\), el perfil
\eqref{eq:perfiles-laterales-fano-semillas-rev3} proporciona exactamente dos
elecciones en \((P_r)_b\) con ese residuo; si no pertenece al soporte, no
proporciona ninguna. Como hay treinta y seis fibras,
\(N_\delta=36\cdot4\cdot2=288\) en el primer caso y cero en el segundo.
Los tres soportes laterales particionan \(R_9^\times\), de modo que una unidad
\(\delta\) determina un único lateral para cada bloque transversal. La lectura
directa de esos soportes da la tabla
\eqref{eq:tabla-theta-delta-fano-semillas-rev3} y agota las \(3!=6\)
biyecciones.

La línea transversal cubre las tres parejas entre \(B_3,B_6,B_9\). Para un
\(B\) fijo, las dos aristas del emparejamiento \(\mu(\theta_\delta(B))\)
cubren una vez los cuatro puntos laterales. Los tres emparejamientos son
distintos y forman la factorización de \(K_4\) en sus tres emparejamientos
perfectos; por tanto cada
pareja lateral aparece exactamente una vez. Las siete ternas contienen así
cada una de las veintiuna parejas de puntos una vez, que es la propiedad
\(\operatorname{S}(2,3,7)\). Recíprocamente, todo plano de Fano que contenga
la línea transversal asigna a cada \(B\) un emparejamiento perfecto distinto
de los otros dos y recupera una de las seis biyecciones anteriores.
\end{proof}

La correspondencia de conjuntos \(\delta\mapsto\theta_\delta\) no es una
acción ni un homomorfismo: \(R_9^\times\simeq C_6\) es abeliano bajo el
producto, mientras \(S_3\) no lo es, y las unidades no forman un subgrupo
aditivo. El desplazamiento \(\delta\) indexa seis traslaciones de los perfiles
del producto.

El mismo cálculo y la misma tabla se verifican en las seis microfibras. Debe
conservarse, no obstante, la diferencia entre soporte y rol: el soporte de
\(f\) coincide con uno de los soportes laterales y, por ello, el observable
\(\rho\) no recupera por sí solo la descomposición \(P\sqcup\{f\}\). Esa
distinción exige la firma aditiva ya declarada en el estado de semilla.

\begin{statusbox}[title={Estatuto y alcance de la incidencia de Fano}]
El \cref{thm:semillas-seis-fano-fase-rev3} es un teorema finito una vez
fijados \(q_\beta\), la orientación y el orden digital: define seis calibres
exactos de incidencia sobre bloques. Organiza fase e incidencia en el dominio
de semillas; no selecciona partícula, masa ni ruta. Su realización física se
efectúa después mediante los morfismos del registro enriquecido.
\end{statusbox}

El paso a la física conserva la cadena
\begin{equation}
 \begin{aligned}
 \text{semilla refinada}
 &\longrightarrow(\text{firmas }+,\times;U_6;w_6;
                   \text{memoria de hoja})\\
 &\longrightarrow\text{registro CT108}
 \longrightarrow(\tau_{12},D_k,\Omega,P_6,\nu_{120},\nu_{270})\\
 &\longrightarrow\text{fibra y multisección}
 \longrightarrow(B_D,B_\Omega)
 \longrightarrow\text{transporte CKM/PMNS}\\
 &\longrightarrow\text{carácter y torres}
 \longrightarrow\text{sección de masa}.
 \end{aligned}
 \label{eq:cadena-semilla-mezcla-masa-autonoma}
\end{equation}
CKM realiza un cambio de base de los operadores quark,
\begin{equation}
 B_d^{(u)}=V_{\rm CKM}B_dV_{\rm CKM}^*,
 \qquad
 R_{\rm act}^{B_d^{(u)}}
 =V_{\rm CKM}R_{\rm act}^{B_d}V_{\rm CKM}^*,
 \label{eq:ckm-transporte-semillas-autonoma}
\end{equation}
y preserva espectro y determinante. La mezcla especifica la representación
del operador en una base de sabor; el registro conserva los autovalores y su
genealogía.

\subsection{Lector neutrínico contenido en el registro}

El sector neutro dispone de un extractor entero construido exclusivamente
con datos del registro. Para cada bloque dodecafásico \(E_m\), definimos
\begin{equation}
 T(m)=\sum_{t\in E_m}(-1)^{\kappa(t)}\varepsilon(t)
       \mathbf1_{\{d(t)=9\}},
 \qquad
 s_C^{(m)}(\nu)=T(m)-T(m+6),
 \label{eq:lector-neutrino-ledger-autonoma}
\end{equation}
y
\begin{equation}
 \tau_m^\nu=\operatorname{sgn}s_C^{(m)}(\nu).
 \label{eq:trit-neutrino-ledger-autonoma}
\end{equation}
Los cuatro triángulos de salto cuatro producen
\begin{equation}
 \nu_{120}=\sum_{j=1}^{4}
 \operatorname{sgn}\!\left(\sum_{m\in\mathcal T_j}\tau_m^\nu\right),
 \qquad
 \nu_{270}=\sum_{m=1}^{12}\tau_m^\nu\tau_{m+3}^\nu.
 \label{eq:invariantes-neutrino-ledger-autonoma}
\end{equation}
La inversión bidireccional transforma
\(\nu_{120}\mapsto-\nu_{120}\) y conserva \(\nu_{270}\). Esta paridad
define la arquitectura histórica de un lector independiente de magnitudes
metrológicas empleadas como blanco. Su aplicación
literal al registro vivo \(81\times108\), tomando como
\(\varepsilon(t)\) el único bit global de signo conservado por esa
proyección, produce
\[
 s_C^{(m)}(\nu)=0,
 \qquad \tau_m^\nu=\nu_{120}=\nu_{270}=0
\]
en las ochenta y una posiciones. Los eventos antipodales han quedado
identificados por la proyección del registro. El resultado local se conserva
como lector recuperado y, a la vez, como test de pérdida de información: el
extractor neutro de
\cref{sec:extractor-neutro-z0-torsion-rev2} actúa antes de esa
identificación y lee el bit bidireccional local, la hoja, la torsión y la
corona antipodal. El cero obtenido después de proyectar no demueve esa
construcción; certifica exactamente qué información borra el lector global.
Los marcos
\(F_\ell,F_\nu\), el lector~\eqref{eq:invariantes-neutrino-ledger-autonoma}
y la forma agregada
\begin{equation}
 G_{ab}=\frac1{\sqrt{|L_a||N_b|}}
 \sum_{c\in L_a}\sum_{d\in N_b}K_{\rm registro}(c,d),
 \qquad U_{\rm int}=\operatorname{polar}(G),
 \label{eq:kernel-ledger-pmns-autonoma}
\end{equation}
fijan así los dos extremos ya construidos del puente PMNS.

La igualdad cardinal entre los cincuenta y seis vértices multiplicativos de
semillas y las cincuenta y seis familias CT108 ha sido sometida a un primer
control independiente de magnitudes metrológicas empleadas como blanco. El
mapa de Hadamard natural sobre una sola semilla alcanza
sólo veintinueve familias; al recorrer las simetrías diédricas, signos y
traslaciones naturales, la cobertura máxima es treinta y cinco. Por tanto el
entrelazador actúa sobre el par de semillas enriquecidas y conserva hoja,
acarreo y memoria. La matriz de incidencia que compare emisiones compartidas
en ese dominio doble deberá tener rango cincuenta y seis, ser natural bajo
orientación y no factorizar por contadores globales. Este criterio fija el
dominio exacto de la siguiente construcción.
